\chapter{Greetings, Farewells, Letters, and Emails}

The appropriate greeting depends on the relationship, situation, time of day,
and region. In unfamiliar, official, or professional situations, begin with
\textit{Sie} unless the other person or the organisation clearly uses
\textit{du}. It is normally safe to match the level of formality used by the
other person.

\section{Everyday spoken greetings}

\begin{center}
\small
\rowcolors{2}{referenceBlueBg}{white}
\begin{tabularx}{\textwidth}{@{}>{\bfseries}l l l X@{}}
\toprule
German & English & Register & Typical use \\
\midrule
Guten Morgen! & Good morning! & neutral/formal & morning \\
Guten Tag! & Good day/hello! & neutral/formal & daytime \\
Guten Abend! & Good evening! & neutral/formal & evening \\
Hallo! & Hello! & neutral/informal & very widely used \\
Hi! & Hi! & informal & friends and familiar contacts \\
Grüß dich! & Hello! & informal & one person, especially regionally \\
Grüß euch! & Hello! & informal & several people, especially regionally \\
Herzlich willkommen! & A warm welcome! & all registers & welcoming a visitor \\
\bottomrule
\end{tabularx}
\end{center}

German normally uses \textit{Guten Tag}, not a literal translation of
\enquote{good afternoon}. \textit{Gute Nacht} is normally said when leaving late or
going to bed, rather than when first greeting somebody in the evening.

\subsection{Regional greetings}

\begin{center}
\small
\rowcolors{2}{referenceGreenBg}{white}
\begin{tabularx}{\textwidth}{@{}>{\bfseries}l l X@{}}
\toprule
Expression & Region/register & Note \\
\midrule
Moin! & northern Germany, informal/neutral & commonly used throughout the day \\
Grüß Gott! & southern Germany and Austria, polite & literally \enquote{greet God} \\
Servus! & southern Germany and Austria, informal & used for hello and goodbye \\
Mahlzeit! & workplace, informal & greeting around lunchtime \\
\bottomrule
\end{tabularx}
\end{center}

Regional expressions are useful to recognise, but the neutral forms
\textit{Guten Tag} and \textit{Hallo} work more widely.

\section{Introductions and small talk}

\begin{center}
\small
\rowcolors{2}{referenceYellowBg}{white}
\begin{tabularx}{\textwidth}{@{}>{\bfseries}X X@{}}
\toprule
German & English \\
\midrule
Wie geht es dir? & How are you? (informal) \\
Wie geht es Ihnen? & How are you? (formal) \\
Danke, gut. Und dir/Ihnen? & Fine, thank you. And you? \\
Darf ich mich vorstellen? & May I introduce myself? \\
Ich heiße \ldots{} / Mein Name ist \ldots{} & My name is \ldots{} \\
Freut mich, dich/Sie kennenzulernen. & Pleased to meet you. \\
Schön, dich/Sie wiederzusehen. & Nice to see you again. \\
Wie war Ihre/deine Reise? & How was your journey? \\
\bottomrule
\end{tabularx}
\end{center}

Use \textit{Herr} or \textit{Frau} with the surname:
\textit{Guten Tag, Frau Wagner}. \textit{Fräulein} is outdated and should not
be used as the normal form of address. Retain an academic title when it is
given, for example \textit{Sehr geehrte Frau Dr.~Weber}.

\section{Saying goodbye}

\begin{center}
\small
\rowcolors{2}{referenceRedBg}{white}
\begin{tabularx}{\textwidth}{@{}>{\bfseries}l l X@{}}
\toprule
German & English & Register/use \\
\midrule
Auf Wiedersehen! & Goodbye! & neutral/formal, face to face \\
Tschüss! & Bye! & neutral/informal \\
Bis bald! & See you soon! & neutral/informal \\
Bis später! & See you later! & neutral/informal \\
Bis morgen/Montag! & See you tomorrow/Monday! & neutral/informal \\
Bis dann! & See you then! & informal \\
Mach's gut! / Macht's gut! & Take care! & informal singular/plural \\
Einen schönen Tag noch! & Have a nice day! & polite/neutral \\
Gute Nacht! & Good night! & late evening/bedtime \\
Auf Wiederhören! & Goodbye! & telephone calls \\
\bottomrule
\end{tabularx}
\end{center}

A natural response is \textit{Danke, gleichfalls!} or
\textit{Danke, Ihnen/dir auch!} To suggest using first names and
\textit{du}, ask \textit{Wollen wir uns duzen?} Until that is established,
continue using \textit{Sie} in a formal relationship.

\section{Useful telephone language}

\begin{center}
\small
\rowcolors{2}{referenceBlueBg}{white}
\begin{tabularx}{\textwidth}{@{}>{\bfseries}X X@{}}
\toprule
German & English \\
\midrule
Guten Tag, hier spricht Alex Morgan. & Hello, Alex Morgan speaking. \\
Könnte ich bitte Frau Klein sprechen? & Could I speak to Ms Klein, please? \\
Einen Moment bitte. & One moment, please. \\
Sie ist leider gerade nicht erreichbar. & Unfortunately, she is unavailable. \\
Kann ich etwas ausrichten? & Can I take a message? \\
Ich rufe später noch einmal an. & I will call again later. \\
Vielen Dank für Ihren Anruf. & Thank you for your call. \\
Auf Wiederhören! & Goodbye! \\
\bottomrule
\end{tabularx}
\end{center}

\section{Choosing a written register}

\begin{center}
\small
\rowcolors{2}{referenceGreenBg}{white}
\begin{tabularx}{\textwidth}{@{}>{\bfseries}l X X@{}}
\toprule
Register & Opening & Closing \\
\midrule
Formal & Sehr geehrte Frau Müller, & Mit freundlichen Grüßen \\
Formal, unknown recipient(s) & Sehr geehrte Damen und Herren, & Mit freundlichen Grüßen \\
Established professional contact & Liebe Frau Müller, & Beste/Freundliche Grüße \\
Neutral/informal & Hallo Anna, & Viele Grüße \\
Personal & Liebe Anna, / Lieber Paul, & Herzliche/Liebe Grüße \\
\bottomrule
\end{tabularx}
\end{center}

\textit{Hallo} is common in ordinary emails, but
\textit{Sehr geehrte Frau + surname} or
\textit{Sehr geehrter Herr + surname} remains the safest opening for a
first formal contact, an application, or communication with an authority.

\section{Structure of a formal letter or email}

\begin{enumerate}[leftmargin=*]
  \item Give the email a short, informative subject. A modern subject line
        does not need the word \textit{Betreff} and normally has no final full
        stop.
  \item Use the person's name when it is known. Put a comma after the
        salutation.
  \item Begin the following line with a lowercase letter unless its first word
        must otherwise be capitalised: \textit{Sehr geehrte Frau Klein,
        vielen Dank \ldots{}}
  \item State the reason for writing early, then divide longer text into clear
        paragraphs.
  \item End with an appropriate closing sentence and a separate complimentary
        close.
  \item Do not put a comma or full stop after a standalone closing such as
        \textit{Mit freundlichen Grüßen}.
  \item Add your full name. A professional email may also include a signature
        block with contact details.
\end{enumerate}

In formal correspondence, the pronouns \textit{Sie, Ihnen}, and
\textit{Ihr/Ihre} are always capitalised. Informal \textit{du, dir, dich,
ihr}, and the corresponding possessives are normally lowercase; in personal
letters, emails, and messages they may alternatively be capitalised. Whichever
option is chosen should be used consistently.

For a printed formal letter, also include the sender and recipient addresses,
place and date, subject, and a handwritten signature. List enclosures beneath
the signature if necessary. Exact positioning may follow an employer's
template or the DIN~5008 business-letter standard.

\section{Useful written phrases}

\begin{center}
\small
\rowcolors{2}{referenceYellowBg}{white}
\begin{tabularx}{\textwidth}{@{}>{\bfseries}l X X@{}}
\toprule
Purpose & German & English \\
\midrule
Opening & Vielen Dank für Ihre Nachricht. & Thank you for your message. \\
Reason & Ich schreibe Ihnen, weil \ldots{} & I am writing because \ldots{} \\
Enquiry & Ich möchte mich nach \ldots{} erkundigen. & I would like to ask about \ldots{} \\
Request & Könnten Sie mir bitte mitteilen, ob \ldots{}? & Could you please tell me whether \ldots{}? \\
Attachment & Im Anhang finden Sie \ldots{} & Please find \ldots{} attached. \\
Apology & Bitte entschuldigen Sie die verspätete Antwort. & Please excuse the late reply. \\
Problem & Leider kann ich den Termin nicht wahrnehmen. & Unfortunately, I cannot attend the appointment. \\
Thanks & Vielen Dank im Voraus. & Thank you in advance. \\
Closing & Ich freue mich auf Ihre Antwort. & I look forward to your reply. \\
Offer & Für Rückfragen stehe ich gern zur Verfügung. & I am happy to answer any questions. \\
\bottomrule
\end{tabularx}
\end{center}

\section{Model emails}

\subsection{Formal email}

\begin{tcolorbox}[
  colback=referenceBlueBg,
  colframe=genderMasculine,
  title={Anfrage zum B1-Abendkurs},
  fonttitle=\bfseries
]
Sehr geehrte Frau Schneider,

\medskip
ich interessiere mich für Ihren Deutschkurs auf Niveau B1. Könnten Sie mir
bitte mitteilen, wann der nächste Abendkurs beginnt und wie viel er kostet?

\medskip
Vielen Dank im Voraus. Ich freue mich auf Ihre Antwort.

\medskip
Mit freundlichen Grüßen

Alex Morgan
\end{tcolorbox}

\subsection{Informal email}

\begin{tcolorbox}[
  colback=referenceRedBg,
  colframe=genderFeminine,
  title={Treffen am Samstag},
  fonttitle=\bfseries
]
Hallo Lena,

\medskip
danke für deine Nachricht. Ich habe am Samstag Zeit. Wollen wir uns um
15~Uhr im Café treffen?

\medskip
Bis bald

Alex
\end{tcolorbox}

\begin{tcolorbox}[
  colback=referenceGreenBg,
  colframe=genderNeuter,
  title={Final check before sending},
  fonttitle=\bfseries
]
Check the recipient and attachment, use one consistent register, include a
clear subject, answer every question or requested point, and make sure the
message has both an opening and an appropriate closing.
\end{tcolorbox}
